\documentclass[11pt]{article}
\usepackage[T1]{fontenc}
\usepackage[utf8]{inputenc}
\usepackage[spanish,es-noquoting]{babel}
\usepackage{amsmath,amssymb,mathtools}
\usepackage[a4paper,margin=2.3cm]{geometry}
\usepackage{graphicx}
\usepackage{tikz}
\usetikzlibrary{arrows.meta,calc,babel}
\usepackage{float}
\usepackage[hidelinks]{hyperref}
\spanishdecimal{.}
\setlength{\parindent}{0pt}
\setlength{\parskip}{0.35em}
\setlength{\abovedisplayskip}{5pt}
\setlength{\belowdisplayskip}{5pt}
\setlength{\abovedisplayshortskip}{4pt}
\setlength{\belowdisplayshortskip}{4pt}
\setlength{\jot}{6pt}
\setlength{\emergencystretch}{2em}
\allowdisplaybreaks[1]
\tikzset{
  >=Stealth,line cap=round,line join=round,
  every node/.style={font=\small},
  eje/.style={-{Stealth[length=1.7mm]},thin},
  vector/.style={-{Stealth[length=2mm]},line width=0.8pt},
  ayuda/.style={densely dashed,thin},
  corriente/.style={-{Stealth[length=2mm]},line width=1pt}
}
\newcommand{\dd}{\,\mathrm{d}}
\newcommand{\ex}{\hat{\mathbf{x}}}
\newcommand{\ey}{\hat{\mathbf{y}}}
\newcommand{\ez}{\hat{\mathbf{z}}}
\newcommand{\ep}{\hat{\boldsymbol{\phi}}}
\newcommand{\er}{\hat{\mathbf r}}
\newcommand{\encabezado}[3]{%
  {\Large\bfseries #1\par}\vspace{2pt}
  {\normalsize #2\hfill #3\par}
  \vspace{6pt}\hrule\vspace{8pt}}
\newcommand{\ideafisica}[1]{%
  \par\medskip\noindent\textit{Interpretación física.} #1\par}
\begin{document}
\encabezado{Magnetostática: problemas 1 a 5}
{Electrodinámica I}{P.S., J.A., J.R.}

\section*{Enunciados de los Problemas}

\subsection*{Problema 1: Biot-Savart}
Considere una espira con forma de polígono regular de $n$ lados que transporta una corriente estacionaria $I$. El polígono se encuentra inscrito en una circunferencia de radio $R$ (es decir, la distancia desde el centro del polígono hasta cualquiera de sus vértices es $R$).
\begin{enumerate}
    \item[a)] Utilice la Ley de Biot-Savart para calcular el campo magnético $\mathbf{B}$ exactamente en el centro geométrico del polígono. 
    \item[b)] Demuestre matemáticamente que, al tomar el límite cuando el número de lados tiende a infinito ($n \to \infty$), el campo magnético obtenido se reduce exactamente a la expresión conocida para el centro de una espira circular de radio $R$.
\end{enumerate}

\subsection*{Problema 2: Ley de Ampère con Corriente Inhomogénea}
Un cilindro conductor infinito de radio $R$ transporta una corriente total $I_{tot}$ dirigida a lo largo de su eje (eje $z$). Sin embargo, la densidad de corriente no es uniforme, sino que aumenta proporcionalmente con la distancia al eje, descrita por la expresión $\mathbf{J}(s) = \alpha s \hat{\mathbf{z}}$, donde $\alpha$ es una constante y $s$ es la distancia radial cilíndrica.
\begin{enumerate}
    \item[a)] Determine la constante $\alpha$ en función de la corriente total $I_{tot}$ y el radio del cilindro $R$.
    \item[b)] Utilizando la Ley de Ampère, calcule el campo magnético $\mathbf{B}(s)$ en la región interior del conductor ($s < R$).
    \item[c)] Calcule el campo magnético $\mathbf{B}(s)$ en la región exterior ($s > R$).
    \item[d)] Grafique de forma cualitativa la magnitud del campo magnético $|\mathbf{B}|$ en función de la distancia $s$, desde $s=0$ hasta $s=3R$, destacando la continuidad en la frontera $s=R$.
\end{enumerate}

\subsection*{Problema 3: Ecuaciones de Frontera y Superposición}
Considere dos láminas planas e infinitas, paralelas al plano $xy$, ubicadas en $z = d$ y $z = -d$. La lámina superior transporta una densidad de corriente superficial uniforme $\mathbf{K}_1 = K_0 \hat{\mathbf{x}}$, mientras que la lámina inferior transporta una corriente opuesta $\mathbf{K}_2 = -K_0 \hat{\mathbf{x}}$.
\begin{enumerate}
    \item[a)] Calcule el vector campo magnético $\mathbf{B}$ en todo el espacio (regiones $z > d$, $-d < z < d$, y $z < -d$) utilizando la Ley de Ampère o el principio de superposición.
    \item[b)] Evalúe analíticamente el salto del campo magnético $\Delta \mathbf{B} = \mathbf{B}_{\text{arriba}} - \mathbf{B}_{\text{abajo}}$ al cruzar exclusivamente la frontera superior en $z = d$. 
    \item[c)] Compruebe de manera explícita que su resultado del inciso anterior satisface la condición de frontera vectorial general: $\Delta \mathbf{B} = \mu_0 (\mathbf{K} \times \hat{\mathbf{n}})$. Identifique correctamente el vector normal unitario $\hat{\mathbf{n}}$.
\end{enumerate}


\subsection*{Problema 4: Potencial Vectorial Magnético del Solenoide}
Un solenoide ideal, infinito, de radio $R$, tiene $n$ vueltas por unidad de longitud y transporta una corriente $I$. Se sabe que el campo magnético en su interior es uniforme ($\mathbf{B} = \mu_0 n I \hat{\mathbf{z}}$) y en su exterior es nulo.
\begin{enumerate}
    \item[a)] Utilizando la forma integral de $\mathbf{B} = \nabla \times \mathbf{A}$ (mediante el Teorema de Stokes), calcule el potencial vectorial magnético $\mathbf{A}(s, \phi, z)$ en las regiones interior ($s < R$) y exterior ($s > R$). Por simetría, asuma que $\mathbf{A}$ apunta en la dirección $\hat{\boldsymbol{\phi}}$ y solo depende de $s$.
    \item[b)] Demuestre explícitamente, calculando el rotacional de los vectores $\mathbf{A}$ obtenidos, que se recuperan los campos magnéticos originales en ambas regiones.
    \item[c)] Compruebe si el potencial vectorial hallado satisface el gauge de Coulomb ($\nabla \cdot \mathbf{A} = 0$).
\end{enumerate}

\subsection*{Problema 5: Expansión Multipolar (Helmholtz vs. Anti-Helmholtz)}
Considere un sistema compuesto por dos espiras circulares idénticas de radio $a$, paralelas al plano $xy$. Sus centros están ubicados en el eje $z$ en $z = d/2$ y $z = -d/2$.
\begin{enumerate}
    \item[a)] \textbf{Caso A (Corrientes en la misma dirección):} Ambas espiras transportan una corriente $I$ en la dirección azimutal $\hat{\boldsymbol{\phi}}$. Calcule el momento dipolar magnético total $\mathbf{m}$ del sistema y escriba la expresión asintótica para el potencial vectorial $\mathbf{A}(\mathbf{r})$ dominante a grandes distancias ($r \gg a, d$).
    \item[b)] \textbf{Caso B (Corrientes en direcciones opuestas):} La espira en $z=d/2$ transporta $I$ en la dirección $\hat{\boldsymbol{\phi}}$, pero la espira en $z=-d/2$ transporta $I$ en $-\hat{\boldsymbol{\phi}}$. Demuestre analíticamente que el momento dipolar magnético $\mathbf{m}$ del sistema es idénticamente cero.
    \item[c)] Concluya físicamente con qué dependencia radial ($1/r^n$) decaerá el potencial vectorial $\mathbf{A}$ y el campo magnético $\mathbf{B}$ dominante a largas distancias para el Caso B.
\end{enumerate}

\clearpage
\section*{Soluciones}

\subsection*{Solución Problema 1: Polígono de $n$ lados}

\paragraph{a) Campo en el centro.}
Tomamos $n\in\mathbb N$, con $n\geq3$, y orientamos la corriente
en sentido antihorario al mirar desde $+z$. Así, las contribuciones
de los lados apuntan en $+\ez$. Si invertimos la corriente, se
invierte el campo.

Primero calculamos la contribución de un segmento rectilíneo de
longitud $L$. Ubicamos el origen en el centro del polígono y el
segmento en $y=-r$, con corriente en $+x$. Aquí $r$ es la distancia
perpendicular al lado, no el radio $R$ de la circunferencia.

\begin{figure}[H]
\centering
\begin{tikzpicture}
  \begin{scope}
    \node at (0,2.6) {Geometría del polígono};
    \draw[ayuda] (0,0) circle (2);
    \draw[thick] (0:2)--(60:2)--(120:2)--(180:2)--(240:2)--(300:2)--cycle;
    \draw[vector] (0,0)--(300:2) node[pos=0.67,right=2pt] {$R$};
    \draw[thin] (0,0)--(0,{-sqrt(3)});
    \node[left] at (0,-0.95) {$r$};
    \draw[thin] (0,{-sqrt(3)+0.18})--(0.18,{-sqrt(3)+0.18})--(0.18,{-sqrt(3)});
    \draw[thin] (0,-0.65) arc[start angle=-90,end angle=-60,radius=0.65];
    \node at (-74:0.98) {$\pi/n$};
    \draw[corriente] (-0.65,{-sqrt(3)})--(0.65,{-sqrt(3)})
      node[midway,below=4pt] {$I$};
    \fill (0,0) circle (1.5pt);
    \node[above left] at (0,0) {$O$};
    \node at (0,0.65) {$\mathbf B_n\parallel+\ez$};
  \end{scope}
  \begin{scope}[xshift=7.1cm]
    \node at (0,2.6) {Contribución de un lado};
    \draw[eje] (-0.35,0)--(2.7,0) node[right] {$x$};
    \draw[eje] (0,-0.3)--(0,1.35) node[above] {$y$};
    \draw[thick] (-2.2,-1.5)--(2.2,-1.5);
    \node[below=4pt] at (-.72,-1.5) {$\mathcal C_1$};
    \node[below] at (-1.65,-1.5) {$y=-r$};
    \coordinate (P) at ({1.5*tan(30)},-1.5);
    \draw[vector] (P)--(0,0) node[pos=0.55,right=3pt] {$\mathbf u$};
    \draw[corriente] (P)--++(0.7,0) node[midway,below=4pt] {$\dd\mathbf l$};
    \draw[ayuda] (0,0)--(0,-1.5);
    \draw[thin] (-0.24,0)--(-0.24,-1.5);
    \draw[thin] (-0.34,0)--(-0.14,0);
    \draw[thin] (-0.34,-1.5)--(-0.14,-1.5);
    \node[left] at (-0.24,-0.75) {$r$};
    \draw[thin] (0,-0.60) arc[start angle=-90,end angle=-60,radius=0.60];
    \node at (-74:0.92) {$\theta$};
    \fill (P) circle (1.5pt);
    \node[above right=2pt] at (P) {$(x,-r)$};
    \fill (0,0) circle (1.5pt);
    \node[above left] at (0,0) {$O$};
    \node[above left=2pt] at (-2.2,-1.5) {$-L/2$};
    \node[above right=2pt] at (2.2,-1.5) {$L/2$};
  \end{scope}
\end{tikzpicture}
\caption{A la izquierda se ilustra el caso $n=6$; el cálculo vale para
cualquier $n\geq3$. A la derecha, $\mathbf u$ va desde el elemento de
corriente al punto de observación $O$.}
\label{fig:poligono}
\end{figure}

La ley de Biot-Savart da el aporte de cada elemento de corriente.
Integramos ese aporte sobre el lado $\mathcal C_1$:
\begin{equation}
  \dd\mathbf B
    =\frac{\mu_0 I}{4\pi}
      \frac{\dd\mathbf l\times\hat{\mathbf u}}{u^2},
  \qquad \mathbf B_1=\int_{\mathcal C_1}\dd\mathbf B.
  \label{p1:biot}
\end{equation}
Con la orientación de la figura~\ref{fig:poligono},
\[
  \dd\mathbf l=\dd x\,\ex,\qquad
  \mathbf u=-x\ex+r\ey,\qquad
  u=\sqrt{x^2+r^2}.
\]
El producto vectorial se calcula como
\begin{align*}
  \dd\mathbf l\times\hat{\mathbf u}
    &=\dd x\,\ex\times\frac{-x\ex+r\ey}{u}\\
    &=\frac{\dd x}{u}
      \Bigl[-x\underbrace{(\ex\times\ex)}_{\mathbf0}
             +r\underbrace{(\ex\times\ey)}_{\ez}\Bigr]\\
    &=\frac{r\dd x}{u}\,\ez.
\end{align*}

Parametrizamos el segmento mediante el ángulo $\theta$ medido
desde la perpendicular:
\begin{equation}
  x=r\tan\theta,\qquad
  \dd x=r\sec^2\theta\dd\theta,\qquad
  u=\frac{r}{\cos\theta}.
  \label{p1:cambio}
\end{equation}
Los extremos del lado corresponden a $\theta=\pm\theta_0$.
Además, $r/u=\cos\theta$, por lo que el producto anterior es
$\cos\theta\dd x\,\ez$.

Sustituyendo \eqref{p1:cambio} en \eqref{p1:biot} e integrando
sobre un lado,
\begin{align*}
  \mathbf B_1
    &=\frac{\mu_0 I}{4\pi}
      \int_{-\theta_0}^{\theta_0}
      \frac{(r\sec^2\theta)\cos\theta}
           {r^2/\cos^2\theta}\dd\theta\,\ez\\
    &=\frac{\mu_0 I}{4\pi r}
      \int_{-\theta_0}^{\theta_0}\cos\theta\dd\theta\,\ez\\
    &=\frac{\mu_0 I}{4\pi r}
      \Bigl[\sin\theta\Bigr]_{-\theta_0}^{\theta_0}\ez\\
    &=\frac{\mu_0 I}{4\pi r}
      \Bigl[\sin\theta_0-\sin(-\theta_0)\Bigr]\ez\\
    &=\frac{\mu_0 I}{4\pi r}(2\sin\theta_0)\ez.
\end{align*}
Por tanto,
\begin{equation}
  \mathbf B_1=\frac{\mu_0 I}{2\pi r}\sin\theta_0\,\ez.
  \label{p1:lado}
\end{equation}

Ahora conectamos este resultado con la geometría del polígono.
El triángulo rectángulo de la figura~\ref{fig:poligono} tiene
hipotenusa $R$ y ángulo $\theta_0=\pi/n$, de modo que
\begin{equation}
  \theta_0=\frac{\pi}{n},
  \qquad
  r=R\cos\left(\frac{\pi}{n}\right).
  \label{p1:geometria}
\end{equation}
Al sustituir \eqref{p1:geometria} en \eqref{p1:lado},
\begin{align*}
  \mathbf B_1
    &=\frac{\mu_0 I}{2\pi R\cos(\pi/n)}
      \sin\left(\frac{\pi}{n}\right)\ez\\
    &=\frac{\mu_0 I}{2\pi R}
      \tan\left(\frac{\pi}{n}\right)\ez.
\end{align*}
Los $n$ lados aportan la misma magnitud y el mismo sentido.
Al sumar,
\begin{equation}
  \boxed{
    \mathbf B_n=n\mathbf B_1
      =\frac{\mu_0 nI}{2\pi R}
       \tan\left(\frac{\pi}{n}\right)\ez.
  }
  \label{p1:poligono}
\end{equation}

\paragraph{b) Límite circular ($n\to\infty$).}
Mantenemos $R$ e $I$ constantes y hacemos el cambio
$\xi=\pi/n$. Cuando $n\to\infty$, $\xi\to0^+$ y $n=\pi/\xi$.
Así, \eqref{p1:poligono} se escribe como
\begin{align*}
  \mathbf B_n
    &=\frac{\mu_0 I}{2\pi R}
      \frac{\pi}{\xi}\tan\xi\,\ez\\
    &=\frac{\mu_0 I}{2R}
      \frac{\tan\xi}{\xi}\,\ez.
\end{align*}
El límite trigonométrico es
\[
  \lim_{\xi\to0^+}\frac{\tan\xi}{\xi}
    =\left(\lim_{\xi\to0^+}\frac{\sin\xi}{\xi}\right)
      \left(\lim_{\xi\to0^+}\frac{1}{\cos\xi}\right)
    =1\cdot1=1.
\]
Por tanto,
\[
  \boxed{
    \lim_{n\to\infty}\mathbf B_n
      =\frac{\mu_0 I}{2R}\,\ez.
  }
\]
Al aumentar $n$, los lados se acercan a la circunferencia de radio
$R$ y recuperamos el campo en el centro de una espira circular.

\clearpage
\subsection*{Solución Problema 2: Ampère en Cilindro Inhomogéneo}

\paragraph{a) Determinación de $\alpha$.}
La expresión general para la corriente que atraviesa una superficie es
\begin{equation}
  I=\int_S\mathbf J\cdot\dd\mathbf a.
  \label{mag:corriente}
\end{equation}
En este problema, elegimos la sección transversal del cilindro con
normal $+\ez$. Los elementos que sustituiremos en
\eqref{mag:corriente} son
\[
  \mathbf J(s')=\alpha s'\ez,
  \qquad
  \dd\mathbf a=s'\dd s'\dd\phi\,\ez.
\]
Ambos vectores son paralelos, como muestra la
figura~\ref{fig:elemento_corriente}. Usamos $s'$ para el radio de
integración, reservando $s$ para la curva amperiana.

\begin{figure}[H]
\centering
\begin{tikzpicture}[line join=round]
  \begin{scope}
    \node at (0,2.65) {Sección transversal $S$};
    \begin{scope}[x={(0.90cm,-0.30cm)},
                  y={(0.55cm,0.48cm)},z={(0cm,1cm)}]
      \path[draw,thick,fill=black!5]
        plot[domain=0:360,samples=101,variable=\t]
          ({1.75*cos(\t)},{1.75*sin(\t)},0) -- cycle;
      % Pequeña celda de área en el plano perpendicular a la corriente.
      \path[draw,fill=black!30]
        plot[domain=25:47,samples=16,variable=\t]
          ({1.25*cos(\t)},{1.25*sin(\t)},0)
        -- plot[domain=47:25,samples=16,variable=\t]
          ({1.02*cos(\t)},{1.02*sin(\t)},0) -- cycle;
      \coordinate (Q) at ({1.135*cos(36)},{1.135*sin(36)},0);
      \draw[vector] (0,0,0)--(1.75,0,0)
        node[pos=0.62,below=2pt] {$R$};
      \draw[eje] (0,0,0)--(0,0,1.98) node[above] {$z$};
    \end{scope}
    \draw[vector] ($(Q)+(-0.12,0)$)--++(0,0.77)
      node[above left=1pt] {$\dd\mathbf a$};
    \draw[corriente] ($(Q)+(0.13,0)$)--++(0,1.14)
      node[above right=1pt] {$\mathbf J$};
    \fill (0,0) circle (1.4pt);
    \node[below left] at (0,0) {$O$};
    \node at (0,-1.55) {$\hat{\mathbf n}=\ez$};
  \end{scope}
  \begin{scope}[xshift=6.3cm]
    \node at (0,2.65) {Elemento de área visto desde $+z$};
    \fill[black!10,even odd rule]
      (0,0) circle (1.25) (0,0) circle (1.02);
    \draw[thick] (0,0) circle (1.75);
    \draw[ayuda] (0,0) circle (1.02);
    \draw[ayuda] (0,0) circle (1.25);
    \path[fill=black!35,draw]
      (25:1.02)--(25:1.25)
      arc[start angle=25,end angle=47,radius=1.25]
      --(47:1.02)
      arc[start angle=47,end angle=25,radius=1.02]--cycle;
    \draw[ayuda] (0,0)--(25:1.47);
    \draw[ayuda] (0,0)--(47:1.47);
    \draw[thin] (25:0.48)
      arc[start angle=25,end angle=47,radius=0.48];
    \node at (36:0.72) {$\dd\phi$};
    \draw[vector] (0,0)--(-105:1.02)
      node[pos=0.65,left] {$s'$};
    \draw[thin] (-55:1.02)--(-55:1.25);
    \draw[thin] (-55:1.13)--(1.72,-1.23)
      node[right] {$\dd s'$};
    \draw[thin] (36:1.25)--(1.92,1.05)
      node[right] {$s'\dd\phi$};
    \fill (0,0) circle (1.4pt);
    \node[above left] at (0,0) {$O$};
    \node at (0,-2.05) {$\dd a=s'\dd s'\dd\phi$};
  \end{scope}
\end{tikzpicture}
\caption{La densidad de corriente atraviesa el elemento de área
en la dirección de su normal. A la derecha, la celda tiene lados
$\dd s'$ y $s'\dd\phi$; al recorrer todo el ángulo forma un anillo
de área $2\pi s'\dd s'$.}
\label{fig:elemento_corriente}
\end{figure}

Sustituyendo en \eqref{mag:corriente} e integrando toda la sección,
\begin{align*}
  I_{\mathrm{tot}}
    &=\int_0^{2\pi}\int_0^R
      (\alpha s'\ez)\cdot(s'\dd s'\dd\phi\,\ez)\\
    &=2\pi\alpha\int_0^R s'^2\dd s'
      \qquad(\ez\cdot\ez=1)\\
    &=2\pi\alpha\Biggl[\frac{s'^3}{3}\Biggr]_0^R
      =\frac{2\pi\alpha R^3}{3}.
\end{align*}
Despejando,
\begin{equation}
  \boxed{\alpha=\frac{3I_{\mathrm{tot}}}{2\pi R^3}.}
  \label{p2:alpha}
\end{equation}

\begin{figure}[H]
\centering
\begin{tikzpicture}
  \begin{scope}
    \draw[fill=black!5] (0,0) circle (1.55);
    \fill[black!15] (0,0) circle (0.95);
    \draw[ayuda] (0,0) circle (0.95);
    \draw[ayuda] (0,0) circle (2.05);
    \draw[vector] (2.05,0)--(2.05,0.72) node[above] {$\mathbf B$};
    \draw[vector] (-65:0.95) arc[start angle=-65,end angle=-20,radius=0.95];
    \draw[vector] (120:2.05)
      arc[start angle=120,end angle=155,radius=2.05]
      node[midway,above left=2pt] {$\dd\mathbf l$};
    \draw[vector] (0,0)--(-40:1.55) node[pos=0.72,below left] {$R$};
    \node[font=\normalsize] at (135:0.65) {$\odot$};
    \node[font=\normalsize] at (35:1.15) {$\odot$};
    \node[font=\normalsize] at (210:1.1) {$\odot$};
    \fill (0,0) circle (1.1pt);
    \node[above right=1pt] at (0,0) {$O$};
    \node at (-0.38,-0.57) {$S_s$};
    \draw[thin] (0,0)--(-100:0.95)
      node[pos=0.68,right=1pt] {$s$};
    \node at (0,2.38) {$\mathcal C_{\mathrm{ext}}$};
    \node[anchor=east] at (-0.70,0.15) {$\mathcal C_s$};
    \node at (0,-2.65) {$\odot:\ \mathbf J\text{ en }+z$};
  \end{scope}
  \begin{scope}[xshift=5.4cm,yshift=-1.4cm,x=1.55cm,y=1.85cm]
    \draw[eje] (0,0)--(3.35,0) node[right] {$s/R$};
    \draw[eje] (0,0)--(0,1.48) node[above] {$B/B_R$};
    \draw[ayuda] (1,0)--(1,1)--(0,1);
    \draw[thick,domain=0:1,samples=45] plot (\x,{\x*\x});
    \draw[thick,domain=1:3,samples=65] plot (\x,{1/\x});
    \fill (1,1) circle (1.5pt);
    \foreach \x in {1,2,3} {
      \draw (\x,0.035)--(\x,-0.035) node[below] {$\x$};
    }
    \node[left] at (0,0) {$0$};
    \node[left] at (0,1) {$1$};
    \node[anchor=east] at (0.86,0.57) {$s^2$};
    \node[above] at (2.3,0.49) {$1/s$};
  \end{scope}
\end{tikzpicture}
\caption{Corte transversal y dos elecciones de curva amperiana.
El disco más oscuro $S_s$ es atravesado por $I_{\mathrm{enc}}(s)$;
$\mathbf B$ y $\dd\mathbf l$ son tangentes a las circunferencias.
Los símbolos $\odot$ indican corriente saliendo del plano.
A la derecha se representa el campo normalizado por
$B_R=B(R)$: crece dentro y decrece fuera, sin saltar en $s=R$.}
\label{fig:cilindro}
\end{figure}

\paragraph{b) Campo interior ($s<R$).}
Por simetría cilíndrica y la regla de la mano derecha,
$\mathbf B=B_\phi(s)\ep$. Elegimos una circunferencia de radio
$s$, como en la figura~\ref{fig:cilindro}. Partimos de la Ley de Ampère,
\begin{equation}
  \oint_{\mathcal C}\mathbf B\cdot\dd\mathbf l
    =\mu_0I_{\mathrm{enc}}.
  \label{mag:ampere}
\end{equation}
En esta curva, $\dd\mathbf l=s\dd\phi\,\ep$ y $B_\phi(s)$ es
constante. Al sustituir estos elementos en \eqref{mag:ampere},
\begin{equation}
  \int_0^{2\pi}(B_\phi(s)\ep)\cdot(s\dd\phi\,\ep)
    =2\pi sB_\phi(s)=\mu_0I_{\mathrm{enc}}(s).
  \label{p2:ampere}
\end{equation}
Para $0<s<R$, la corriente encerrada se obtiene integrando solo
hasta $s$. Usando \eqref{p2:alpha},
\begin{equation}
  I_{\mathrm{enc}}(s)
    =2\pi\alpha\int_0^s s'^2\dd s'
    =\frac{2\pi\alpha s^3}{3}
    =I_{\mathrm{tot}}\left(\frac{s}{R}\right)^3.
  \label{p2:Ienc}
\end{equation}
Sustituyendo \eqref{p2:Ienc} en \eqref{p2:ampere} y despejando,
\begin{align*}
  B_\phi(s)
    &=\frac{\mu_0I_{\mathrm{enc}}(s)}{2\pi s}
      =\frac{\mu_0I_{\mathrm{tot}}}{2\pi s}\frac{s^3}{R^3}\\
    &=\frac{\mu_0I_{\mathrm{tot}}}{2\pi R^3}s^2.
\end{align*}

\paragraph{c) Campo exterior ($s>R$).}
La curva encierra toda la sección del conductor, así que
$I_{\mathrm{enc}}=I_{\mathrm{tot}}$. Usando \eqref{p2:ampere},
\[
  2\pi sB_\phi(s)=\mu_0I_{\mathrm{tot}}
  \quad\Longrightarrow\quad
  B_\phi(s)=\frac{\mu_0I_{\mathrm{tot}}}{2\pi s}.
\]
Reuniendo las dos regiones,
\begin{equation}
  \boxed{
    \mathbf B(s)=
    \begin{cases}
      \displaystyle
      \frac{\mu_0I_{\mathrm{tot}}}{2\pi R^3}s^2\ep,
        &0<s<R,\\[7pt]
      \displaystyle
      \frac{\mu_0I_{\mathrm{tot}}}{2\pi s}\ep,
        &s>R.
    \end{cases}
  }
  \label{p2:B}
\end{equation}
En el eje, $\mathbf B(0)=\mathbf0$.

\paragraph{d) Gráfico cualitativo y continuidad.}
Al evaluar \eqref{p2:B} a ambos lados de $R$, obtenemos
\[
  B(R^-)
    =\frac{\mu_0I_{\mathrm{tot}}}{2\pi R^3}R^2
    =\frac{\mu_0I_{\mathrm{tot}}}{2\pi R}
    =B(R^+).
\]
Definimos $B_R:=B(R)=\mu_0I_{\mathrm{tot}}/(2\pi R)$.
Para trazar el gráfico de la figura~\ref{fig:cilindro} hasta $3R$,
usamos
\[
  \frac{B(s)}{B_R}
    =\begin{cases}
       (s/R)^2,&0\leq s/R\leq1,\\[3pt]
       R/s,&1\leq s/R\leq3.
     \end{cases}
\]
\ideafisica{Dentro del conductor, el campo depende de cómo se
distribuye la corriente y crece como $s^2$ desde cero. Fuera,
depende de la corriente total y decrece como $1/s$. Alcanza su
máximo en $R$ y es continuo allí, aunque su pendiente cambia.}

\par\bigskip\hrule\bigskip
\subsection*{Solución Problema 3: Superposición y Fronteras}

\paragraph{a) Campo de una lámina individual y superposición.}
Primero resolvemos una sola lámina en $z=0$, con
$\mathbf K=K_0\ex$, y después sumaremos los campos de las dos
láminas del enunciado. Podemos pensar la lámina como muchos alambres
paralelos al eje $x$: sus contribuciones normales a la lámina se
cancelan por pares y queda un campo en la dirección $y$.
Con el pulgar hacia $+x$, la regla de la mano derecha da
\textbf{$-\ey$ arriba y $+\ey$ abajo}.
Por simetría traslacional en $x$ e $y$, \textbf{el campo solo depende
de $z$}, y a alturas opuestas $z=\pm h$ tiene igual magnitud $B$.

\begin{figure}[H]
\centering
\begin{tikzpicture}[scale=1.17,every node/.append style={font=\normalsize}]
  \fill[black!8] (-2.40,-1.30) rectangle (2.40,1.30);
  \draw[eje] (-3.85,-1.65)--(-3.85,1.72) node[above] {$z$};
  \draw[eje] (-3.85,-1.65)--(-3.32,-1.65) node[right] {$y$};
  \draw[line width=.85pt] (-3.22,0)--(3.22,0);
  \foreach \x in {-2.95,-1.60,0,1.60,2.95}{
    \node[fill=white,inner sep=1pt] at (\x,0) {$\odot$};
  }
  \node[right=3pt] at (3.23,0) {$\mathbf K=K_0\ex$};

  \draw[ayuda] (-2.40,-1.30) rectangle (2.40,1.30);
  \draw[vector] (-2.40,.98)--(-2.40,.38)
    node[midway,left=3pt] {$\mathcal C_1$};
  \draw[vector] (.75,1.30)--(-.75,1.30)
    node[midway,above=2pt] {$\mathcal C_2$};
  \draw[vector] (2.40,-.98)--(2.40,-.38)
    node[midway,right=3pt] {$\mathcal C_3$};
  \draw[vector] (-.75,-1.30)--(.75,-1.30)
    node[midway,below=2pt] {$\mathcal C_4$};

  \draw[vector] (.85,.76)--(-.85,.76)
    node[midway,below=2pt] {$\mathbf B=-B\ey$};
  \draw[vector] (-.85,-.76)--(.85,-.76)
    node[midway,above=2pt] {$\mathbf B=+B\ey$};
  \node[anchor=west] at (-3.20,1.48) {$\hat{\mathbf n}_S=\ex$};
  \node[right=3pt] at (2.40,1.30) {$z=h$};
  \node[right=3pt] at (2.40,-1.30) {$z=-h$};
  \draw[thin] (-2.40,-1.60)--(-2.40,-1.87)
    (2.40,-1.60)--(2.40,-1.87);
  \draw[thin,<->] (-2.40,-1.77)--node[below=2pt] {$L$} (2.40,-1.77);
\end{tikzpicture}
\caption{Corte $yz$ de una lámina con corriente hacia $+x$,
saliendo del plano. El rectángulo es la curva amperiana; su
recorrido es antihorario y su superficie tiene normal $+\ex$.
Las flechas del borde indican el sentido de cada tramo $\mathcal C_i$.}
\label{fig:lamina}
\end{figure}

Elegimos el rectángulo de la figura~\ref{fig:lamina}, de anchura
$L$, con lados horizontales en $z=\pm h$. Separamos en cuatro tramos
la circulación de la ley de Ampère \eqref{mag:ampere}. En cada tramo,
$dl$ es la longitud que avanzamos en el sentido de su flecha.
\begin{align*}
  \int_{\mathcal C_1}\mathbf B\cdot\dd\mathbf l
    &=\int_{\mathcal C_1}\mathbf B(z)\cdot(-\ez)\dd l=0,\\
  \int_{\mathcal C_2}\mathbf B\cdot\dd\mathbf l
    &=\int_0^L B\underbrace{(-\ey)\cdot(-\ey)}_{1}\dd l=BL,\\
  \int_{\mathcal C_3}\mathbf B\cdot\dd\mathbf l
    &=\int_{\mathcal C_3}\mathbf B(z)\cdot\ez\dd l=0,\\
  \int_{\mathcal C_4}\mathbf B\cdot\dd\mathbf l
    &=\int_0^L B\underbrace{\ey\cdot\ey}_{1}\dd l=BL .
\end{align*}
El campo es perpendicular al recorrido en $\mathcal C_1$ y
$\mathcal C_3$. Por tanto,
\[
  \oint_{\mathcal C}\mathbf B\cdot\dd\mathbf l
  =\sum_{i=1}^{4}\int_{\mathcal C_i}\mathbf B\cdot\dd\mathbf l
  =0+BL+0+BL=2BL .
\]
Aquí $K_0$ es corriente por unidad de anchura, medida en $\mathrm{A/m}$.
La corriente atraviesa el rectángulo solo donde este corta la lámina,
en un tramo de longitud $L$. Por eso,
\[
  I_{\mathrm{enc}}=\int_0^L K_0\dd l=K_0L.
\]
Sustituyendo ambas expresiones en \eqref{mag:ampere},
\[
  2BL=\mu_0K_0L
  \quad\Longrightarrow\quad
  B=\frac{\mu_0K_0}{2}.
\]
Como $h$ es arbitrario y el resultado no depende de él,
el campo es uniforme en cada semiespacio. Para esta lámina,
\begin{equation}
  \mathbf B_{\text{lámina}}(z)=
  \begin{cases}
    \displaystyle-\frac{\mu_0K_0}{2}\ey,&z>0,\\[5pt]
    \displaystyle+\frac{\mu_0K_0}{2}\ey,&z<0.
  \end{cases}
  \label{p3:lamina}
\end{equation}

Para aplicar \eqref{p3:lamina} a cada lámina, escribimos
$B_0=\mu_0K_0/2$ para la magnitud del campo de una lámina.
La superior, en $z=d$, produce $\mathbf B_1=-B_0\ey$ por
encima y $\mathbf B_1=+B_0\ey$ por debajo. En la inferior,
en $z=-d$, la corriente va en sentido opuesto. Por eso
$\mathbf B_2=+B_0\ey$ por encima y $\mathbf B_2=-B_0\ey$
por debajo.

\begin{figure}[H]
\centering
\begin{tikzpicture}[scale=1.10,every node/.append style={font=\normalsize},
  plano/.style={line width=.72pt},
  referencia/.style={black!62,densely dashed,line width=.45pt},
  campo/.style={vector,line width=.85pt},
  giro/.style={font=\small,inner sep=.8pt,fill=white}
]
  \fill[black!10] (6.65,-1.0) rectangle (9.35,1.0);
  \draw[eje] (-2.20,-1.84)--(-2.20,2.05) node[above] {$z$};
  \draw[eje] (-2.20,-1.84)--(-1.58,-1.84) node[right] {$y$};
  \draw[thin] (-2.29,1)--(-2.11,1) (-2.29,-1)--(-2.11,-1);
  \node[left=3pt] at (-2.20,1) {$d$};
  \node[left=3pt] at (-2.20,-1) {$-d$};

  \begin{scope}
    \node at (0,2.03) {(a) Campo $\mathbf B_1$};
    \draw[plano] (-1.35,1)--(1.35,1);
    \draw[referencia] (-1.35,-1)--(1.35,-1);
    \foreach \x in {-.91,0,.91}
      \node[giro] at (\x,1) {$\odot$};
    \draw[campo] (.56,1.51)--(-.56,1.51);
    \draw[campo] (-.56,0)--(.56,0);
    \draw[campo] (-.56,-1.51)--(.56,-1.51);
  \end{scope}

  \begin{scope}[xshift=4cm]
    \node at (0,2.03) {(b) Campo $\mathbf B_2$};
    \draw[referencia] (-1.35,1)--(1.35,1);
    \draw[plano] (-1.35,-1)--(1.35,-1);
    \foreach \x in {-.91,0,.91}
      \node[giro] at (\x,-1) {$\otimes$};
    \draw[campo] (-.56,1.51)--(.56,1.51);
    \draw[campo] (-.56,0)--(.56,0);
    \draw[campo] (.56,-1.51)--(-.56,-1.51);
  \end{scope}

  \begin{scope}[xshift=8cm]
    \node at (0,2.03) {(c) Campo $\mathbf B_1+\mathbf B_2$};
    \draw[plano] (-1.35,1)--(1.35,1) (-1.35,-1)--(1.35,-1);
    \foreach \x in {-.91,0,.91}{
      \node[giro] at (\x,1) {$\odot$};
      \node[giro] at (\x,-1) {$\otimes$};
    }
    \node at (0,1.51) {$\mathbf 0$};
    \node at (0,-1.51) {$\mathbf 0$};
    \draw[campo] (-1.12,0)--(1.12,0);
    \draw[vector] (1.35,1)--(1.35,1.56)
        node[right=1pt] {$\hat{\mathbf n}=\ez$};
  \end{scope}
\end{tikzpicture}
\caption{Corte $yz$ de las dos láminas. Los puntos indican
corriente hacia $+x$ y las cruces, hacia $-x$. En (a) y (b)
se muestra el campo de cada lámina y se marca con línea
discontinua la posición de la otra. En (c) se suman los campos
de las dos láminas en cada región.}
\label{fig:doslaminas}
\end{figure}

Por superposición,
\begin{equation}
  \mathbf B(z)=\mathbf B_1(z)+\mathbf B_2(z).
  \label{p3:superposicion}
\end{equation}
Sustituimos en \eqref{p3:superposicion} los campos que se ven
en la figura~\ref{fig:doslaminas},
\begin{align*}
  z>d:\quad
  \mathbf B(z)&=(-B_0\ey)+(+B_0\ey)=\mathbf 0,\\[3pt]
  -d<z<d:\quad
  \mathbf B(z)&=(+B_0\ey)+(+B_0\ey)=\mu_0 K_0\ey,\\[3pt]
  z<-d:\quad
  \mathbf B(z)&=(+B_0\ey)+(-B_0\ey)=\mathbf 0.
\end{align*}
Así, el campo neto queda
\begin{equation}
  \boxed{
    \mathbf B(z)=
    \begin{cases}
      \mathbf0,&z>d,\\[3pt]
      \mu_0K_0\ey,&-d<z<d,\\[3pt]
      \mathbf0,&z<-d.
    \end{cases}
  }
  \label{p3:Btotal}
\end{equation}

\paragraph{b) Salto en $z=d$.}
El salto compara el campo a ambos lados de la \emph{misma} lámina
superior. Justo debajo de ella seguimos en la región intermedia.
De \eqref{p3:Btotal},
\[
  \mathbf B(d^+)=\mathbf0,
  \qquad \mathbf B(d^-)=\mu_0K_0\ey.
\]
Sustituyendo en la definición del salto,
\begin{equation}
  \Delta\mathbf B
    =\mathbf B(d^+)-\mathbf B(d^-)
    =\mathbf0-\mu_0K_0\ey
    =-\mu_0K_0\ey.
  \label{p3:salto}
\end{equation}

\paragraph{c) Condición de frontera general.}
Partimos de $\Delta\mathbf B=\mu_0(\mathbf K\times\hat{\mathbf n})$.
Como definimos el salto de abajo hacia arriba, elegimos
$\hat{\mathbf n}=\ez$, como en la figura~\ref{fig:doslaminas}.
Esta es la normal de la \emph{lámina} en el plano $xy$; la normal
$\hat{\mathbf n}_S=\ex$ de antes correspondía al \emph{rectángulo}
en el plano $yz$.

En esta frontera usamos $\mathbf K_1=K_0\ex$. La lámina inferior
no se cruza y su campo es el mismo justo arriba y justo abajo de
$z=d$, por lo que su contribución se cancela en el salto.
Sustituyendo en la condición del enunciado,
\[
  \mu_0(\mathbf K_1\times\hat{\mathbf n})
    =\mu_0K_0(\ex\times\ez)
    =-\mu_0K_0\ey
    =\Delta\mathbf B,
\]
que coincide con \eqref{p3:salto}.
\ideafisica{Entre las láminas, los campos tienen el mismo sentido
y se suman; fuera, se cancelan. Al cruzar la lámina superior,
la corriente superficial produce el salto de la componente
tangencial. La componente normal es cero a ambos lados.}


\clearpage
\subsection*{Solución Problema 4: Potencial Vectorial de Solenoide}

\paragraph{a) Cálculo de $\mathbf A$ mediante Stokes.}
Escribimos $B_0=\mu_0 nI$: el campo es $B_0\ez$ dentro del
solenoide y cero fuera. Como indica el enunciado, buscamos
$\mathbf A=A_\phi(s)\ep$, regular en el eje.
Tomamos una circunferencia $\mathcal C_s$ de radio $s$, recorrida
en sentido antihorario al mirar desde $+z$, y el disco $S_s$
que limita, con normal $+\ez$. El Teorema de Stokes da
\begin{equation}
  \oint_{\mathcal C_s}\mathbf A\cdot\dd\mathbf l
    =\int_{S_s}(\nabla\times\mathbf A)\cdot\dd\mathbf a
    =\int_{S_s}\mathbf B\cdot\dd\mathbf a
    =\Phi_B(s).
  \label{p4:stokes}
\end{equation}

\begin{figure}[H]
\centering
\begin{tikzpicture}
  \begin{scope}
    \node at (0,2.65) {Curva interior: $s<R$};
    \draw[fill=black!4] (0,0) circle (1.35);
    \fill[black!16] (0,0) circle (1.03);
    \draw[ayuda] (0,0) circle (1.03);
    \node at (-1.62,-.33) {$\mathcal C_s$};
    \draw[vector] (80:1.03)
      arc[start angle=80,end angle=110,radius=1.03];
    \node at (95:1.72) {$\dd\mathbf l$};
    \draw[vector] (1.03,0)--(1.03,0.72)
      node[above] {$\mathbf A$};
    \draw[vector] (0,0)--(-135:1.35)
      node[pos=0.70,below right] {$R$};
    \draw[thin] (0,0)--(-40:1.03)
      node[pos=0.67,above right] {$s$};
    \node[font=\normalsize] at (0,0.18) {$\odot$};
    \node at (0,0.55) {$\mathbf B=B_0\ez$};
    \node at (0,-2.48) {$\Phi_B=B_0\pi s^2$};
  \end{scope}
  \begin{scope}[xshift=7.3cm]
    \node at (0,2.65) {Curva exterior: $s>R$};
    \draw[fill=black!16] (0,0) circle (1.35);
    \draw[ayuda] (0,0) circle (2.03);
    \node at (-2.43,-.43) {$\mathcal C_s$};
    \draw[vector] (120:2.03)
      arc[start angle=120,end angle=145,radius=2.03];
    \node at (130:2.35) {$\dd\mathbf l$};
    \draw[vector] (2.03,0)--(2.03,0.72)
      node[above] {$\mathbf A$};
    \draw[vector] (0,0)--(-135:1.35)
      node[pos=0.70,below right] {$R$};
    \draw[thin] (0,0)--(-40:2.03)
      node[pos=0.80,above right] {$s$};
    \node[font=\normalsize] at (0,0.18) {$\odot$};
    \node at (0,0.68) {$\mathbf B=B_0\ez$};
    \node at (0,1.70) {$\mathbf B=\mathbf0$};
    \node at (0,-2.48) {$\Phi_B=B_0\pi R^2$};
  \end{scope}
\end{tikzpicture}
\caption{Cortes $xy$ del solenoide. Las curvas discontinuas son los
recorridos de Stokes y la zona más oscura indica dónde contribuye
$\mathbf B$ al flujo. $\mathbf A$ y $\dd\mathbf l$ son tangentes;
$\mathbf B$ y la normal del disco apuntan hacia $+z$ ($\odot$).}
\label{fig:solenoide}
\end{figure}

Sobre $\mathcal C_s$, $\mathbf A$ es tangente y su magnitud es
constante. Identificamos $\dd\mathbf l=s\dd\phi\,\ep$ y
$\dd\mathbf a=s'\dd s'\dd\phi\,\ez$, con la orientación de la
figura~\ref{fig:solenoide}. Al sustituir en \eqref{p4:stokes},
\[
  \int_0^{2\pi}(A_\phi(s)\ep)\cdot(s\dd\phi\,\ep)
    =2\pi sA_\phi(s)=\Phi_B(s).
\]
Por tanto,
\begin{equation}
  A_\phi(s)=\frac{\Phi_B(s)}{2\pi s},
  \qquad s>0.
  \label{p4:Aflujo}
\end{equation}

\textbf{Interior ($s<R$).}
\textbf{Todo el disco tiene campo uniforme} $B_0\ez$, paralelo a su normal.
El flujo es $\Phi_B=\int_{S_s} B_0\dd a=B_0\pi s^2$.
Sustituyendo en \eqref{p4:Aflujo},
\[
  A_{\phi,\mathrm{int}}(s)
    =\frac{B_0\pi s^2}{2\pi s}
    =\frac{B_0s}{2}.
\]

\textbf{Exterior ($s>R$).}
El campo solo contribuye al flujo hasta $R$. Calculamos ese flujo
y lo sustituimos en \eqref{p4:Aflujo}:
\begin{align*}
  \Phi_B(s)
    &=\int_0^R B_0\,2\pi s'\dd s'
      =B_0\pi R^2,\\
  A_{\phi,\mathrm{ext}}(s)
    &=\frac{B_0\pi R^2}{2\pi s}
      =\frac{B_0R^2}{2s}.
\end{align*}
Al reemplazar $B_0=\mu_0nI$, obtenemos
\begin{equation}
  \boxed{
    \mathbf A(s)=
    \begin{cases}
      \displaystyle\frac{\mu_0 nI}{2}s\,\ep,
        &0<s<R,\\[7pt]
      \displaystyle\frac{\mu_0 nIR^2}{2s}\,\ep,
        &s>R.
    \end{cases}
  }
  \label{p4:A}
\end{equation}
En el eje definimos $\mathbf A(0)=\mathbf0$, y los dos límites
en $R$ coinciden.
\ideafisica{Fuera del solenoide, el flujo es constante porque
el campo solo existe hasta $R$: aumentar el disco no agrega
flujo. Aunque allí $\mathbf B=\mathbf0$, la circulación de
$\mathbf A$ sigue siendo igual al flujo interior y obtenemos
$A_\phi\propto1/s$.}

\paragraph{b) Verificación mediante el rotacional.}
En coordenadas cilíndricas, para el campo azimutal
$\mathbf A=A_\phi(s)\ep$, con $A_s=A_z=0$ y sin dependencia
en $\phi$ ni en $z$, la expresión general del rotacional se reduce a
\begin{equation}
  \nabla\times\mathbf A
    =\frac{1}{s}\frac{\dd(sA_\phi)}{\dd s}\ez.
  \label{p4:rotacional}
\end{equation}
Sustituyendo \eqref{p4:A} en \eqref{p4:rotacional}, en el interior
\begin{align*}
  \nabla\times\mathbf A_{\mathrm{int}}
    &=\frac{1}{s}\frac{\dd}{\dd s}
      \Biggl[\frac{\mu_0 nI}{2}s^2\Biggr]\ez\\
    &=\frac{\mu_0 nI}{2s}(2s)\ez
      =\mu_0 nI\ez
      =\mathbf B_{\mathrm{int}}.
\end{align*}
Para el potencial exterior,
\begin{align*}
  \nabla\times\mathbf A_{\mathrm{ext}}
    &=\frac{1}{s}\frac{\dd}{\dd s}
      \Biggl[s\,\frac{\mu_0 nIR^2}{2s}\Biggr]\ez
      =\frac{1}{s}\frac{\dd}{\dd s}
        \Biggl[\frac{\mu_0 nIR^2}{2}\Biggr]\ez\\
    &=\mathbf0
      =\mathbf B_{\mathrm{ext}}.
\end{align*}
Así recuperamos los campos del enunciado en ambas regiones.

\paragraph{c) Gauge de Coulomb.}
La divergencia de un campo vectorial en coordenadas cilíndricas es
\[
  \nabla\cdot\mathbf A
    =\frac{1}{s}\frac{\partial(sA_s)}{\partial s}
      +\frac{1}{s}\frac{\partial A_\phi}{\partial\phi}
      +\frac{\partial A_z}{\partial z}.
\]
En \eqref{p4:A}, $A_s=A_z=0$ y $A_\phi$ no depende de $\phi$,
de modo que los tres términos son nulos. Así,
\textbf{$\nabla\cdot\mathbf A=0$ y el potencial satisface el gauge
de Coulomb en ambas regiones.}

\clearpage
\subsection*{Solución Problema 5: Expansión Multipolar}

Ubicamos el origen en el punto medio entre las espiras. El punto
de observación tiene vector posición $\mathbf r$, con
$r=|\mathbf r|$, y el ángulo polar $\theta$ se mide desde $+z$.
Trabajamos con $a>0$ y $d>0$, sin fijar una relación entre ambos.

\begin{figure}[H]
\centering
\begin{tikzpicture}
  \foreach \xx/\caso/\sgn in {0/A/1,7.3/B/-1} {
    \begin{scope}[xshift=\xx cm]
      \node at (0,3.00) {Caso \caso};
      \draw[ayuda] (0,-2.2)--(0,2.22);
      \draw[eje] (0,2.12)--(0,2.63) node[right] {$z$};
      \draw[thick] (0,1.02) ellipse[x radius=1.48,y radius=0.35];
      \draw[thick] (0,-1.02) ellipse[x radius=1.48,y radius=0.35];
      \draw[corriente]
        plot[domain=230:310,samples=25]
          ({1.48*cos(\x)},{1.02+0.35*sin(\x)});
      \ifnum\sgn=1
        \draw[corriente]
          plot[domain=230:310,samples=25]
            ({1.48*cos(\x)},{-1.02+0.35*sin(\x)});
        \draw[vector] (0,-1.02)--(0,-0.26)
          node[pos=0.55,right=3pt] {$\mathbf m_2$};
      \else
        \draw[corriente]
          plot[domain=310:230,samples=25]
            ({1.48*cos(\x)},{-1.02+0.35*sin(\x)});
        \draw[vector] (0,-1.02)--(0,-1.96)
          node[pos=0.62,right=3pt] {$\mathbf m_2$};
      \fi
      \draw[vector] (0,1.02)--(0,1.98)
        node[pos=0.62,right=3pt] {$\mathbf m_1$};
      \draw[thin] (0,1.02)--(1.48,1.02)
        node[pos=0.70,above=2pt] {$a$};
      \node[anchor=west] at (1.65,1.02) {$z=d/2$};
      \node[anchor=west] at (1.65,-1.02) {$z=-d/2$};
      \node[left=2pt] at (-1.48,1.02) {$\mathcal C_1$};
      \node[left=2pt] at (-1.48,-1.02) {$\mathcal C_2$};
      \node at (0.95,0.48) {$I$};
      \node at (0.95,-1.57) {$I$};
      \fill (0,0) circle (1.3pt);
      \node[left=3pt] at (0,0) {$O$};
    \end{scope}
  }
\end{tikzpicture}
\caption{Las flechas sobre las espiras indican la corriente.
Los momentos se orientan con la regla de la mano derecha:
en A se suman y en B son opuestos. Los dibujos son esquemáticos.}
\label{fig:espiras}
\end{figure}

\paragraph{a) Caso A: corrientes en la misma dirección.}
El momento dipolar magnético de una espira plana es
\begin{equation}
  \mathbf m_i=I_i\mathbf a_i,
  \qquad
  \mathbf a_i=\pi a^2\ez,
  \label{p5:dipolo}
\end{equation}
donde orientamos el vector de área hacia $+z$ y tomamos
$I_i>0$ para la corriente en $+\ep$.

En este caso, $I_1=I_2=I$. Al sustituir en \eqref{p5:dipolo},
\[
  \mathbf m_1=I\pi a^2\ez,
  \qquad
  \mathbf m_2=I\pi a^2\ez.
\]
Sumando las contribuciones,
\begin{align}
  \mathbf m_{\mathrm{tot}}
    &=\mathbf m_1+\mathbf m_2
      \notag\\
    &=I\pi a^2\ez+I\pi a^2\ez
      \notag\\
    &=\boxed{2I\pi a^2\ez}.
      \label{p5:mA}
\end{align}

Para corrientes localizadas usamos el potencial magnetostático
en el gauge de Coulomb, que se anula al infinito. Su término
dipolar es
\begin{equation}
  \mathbf A_{\mathrm{dip}}(\mathbf r)
    =\frac{\mu_0}{4\pi}
      \frac{\mathbf m_{\mathrm{tot}}\times\er}{r^2}.
  \label{p5:Adip}
\end{equation}
Por la definición de los versores esféricos,
\[
  \ez\times\er
    =\sin\theta(-\sin\phi\,\ex+\cos\phi\,\ey)
    =\sin\theta\,\ep.
\]
Sustituyendo \eqref{p5:mA} en \eqref{p5:Adip},
\begin{align*}
  \mathbf A_{\mathrm A}(\mathbf r)
    &\simeq\frac{\mu_0}{4\pi}
      \frac{2I\pi a^2(\ez\times\er)}{r^2}\\
    &=\frac{2\mu_0I\pi a^2}{4\pi r^2}
      \sin\theta\,\ep\\
    &=\boxed{\frac{\mu_0Ia^2}{2r^2}\sin\theta\,\ep},
      \qquad r\gg a,d.
\end{align*}
A grandes distancias, el sistema se ve como un dipolo de momento
$2I\pi a^2\ez$, como anticipa la figura~\ref{fig:espiras}.

\paragraph{b) Caso B: corrientes en direcciones opuestas.}
Mantenemos la orientación $+\ez$ para los vectores de área.
La inversión de la corriente inferior corresponde a $I_2=-I$.
Usando \eqref{p5:dipolo},
\[
  \mathbf m_1=I\pi a^2\ez,
  \qquad
  \mathbf m_2=(-I)\pi a^2\ez=-I\pi a^2\ez.
\]
Por tanto,
\begin{equation}
  \boxed{
    \mathbf m_{\mathrm{tot}}
      =I\pi a^2\ez-I\pi a^2\ez
      =\mathbf0.
  }
  \label{p5:mB}
\end{equation}
Se anula el momento dipolar total, pero las espiras no están
en el mismo lugar. Por eso esta cancelación no implica que
el campo sea nulo en todo el espacio.

\paragraph{c) Decaimiento dominante en el caso B.}
Para una espira cerrada, el término $1/r$ del potencial se anula
porque $\oint\dd\mathbf l'=\mathbf0$. En el par,
\[
  \mathbf A_{1/r}(\mathbf r)
    =\frac{\mu_0}{4\pi r}
      \sum_{i=1}^2 I_i\oint_{\mathcal C_i}\dd\mathbf l'_i
    =\mathbf0.
\]
Además, \eqref{p5:mB} anula el término dipolar, que decae como
$1/r^2$. Para identificar el primer término que queda,
aprovechamos que las espiras son idénticas y están desplazadas.

Sea $\mathbf A_0(x,y,z)$ el potencial de una espira con corriente
$I$ centrada en el origen. Por superposición,
\begin{equation}
  \mathbf A_{\mathrm B}(x,y,z)
    =\mathbf A_0(x,y,z-d/2)
      -\mathbf A_0(x,y,z+d/2).
  \label{p5:desplazadas}
\end{equation}
A grandes distancias, $d/r\ll1$, y una expansión de Taylor da
\begin{align*}
  \mathbf A_0(x,y,z-d/2)
    &\simeq\mathbf A_0(x,y,z)
      -\frac d2\frac{\partial\mathbf A_0}{\partial z},\\
  \mathbf A_0(x,y,z+d/2)
    &\simeq\mathbf A_0(x,y,z)
      +\frac d2\frac{\partial\mathbf A_0}{\partial z}.
\end{align*}
Sustituyendo en \eqref{p5:desplazadas},
\begin{align}
  \mathbf A_{\mathrm B}
    &\simeq
      \Biggl[\mathbf A_0-\frac d2
        \frac{\partial\mathbf A_0}{\partial z}\Biggr]
      -\Biggl[\mathbf A_0+\frac d2
        \frac{\partial\mathbf A_0}{\partial z}\Biggr]
      \notag\\
    &=-d\,\frac{\partial\mathbf A_0}{\partial z}.
      \label{p5:gradiente}
\end{align}

Ahora usamos el término dipolar de una sola espira,
con $m_0:=I\pi a^2$. De \eqref{p5:Adip},
\[
  \mathbf A_0(\mathbf r)
    \simeq\frac{\mu_0m_0}{4\pi}
      \frac{\ez\times\mathbf r}{r^3}
    =\frac{\mu_0m_0}{4\pi}
      \frac{-y\ex+x\ey}{(x^2+y^2+z^2)^{3/2}}.
\]
El numerador no depende de $z$. Al sustituir en
\eqref{p5:gradiente} y derivar,
\begin{align*}
  \mathbf A_{\mathrm B}
    &\simeq-\frac{\mu_0m_0d}{4\pi}
      (-y\ex+x\ey)
      \frac{\partial}{\partial z}
      (x^2+y^2+z^2)^{-3/2}\\
    &=-\frac{\mu_0m_0d}{4\pi}
      (-y\ex+x\ey)
      \Biggl[-\frac32(x^2+y^2+z^2)^{-5/2}(2z)\Biggr]\\
    &=\frac{3\mu_0m_0d}{4\pi}
      \frac{z(-y\ex+x\ey)}{r^5}.
\end{align*}
Como $z=r\cos\theta$ y $-y\ex+x\ey=r\sin\theta\,\ep$,
\begin{align}
  \mathbf A_{\mathrm B}
    &\simeq\frac{3\mu_0m_0d}{4\pi}
      \frac{r^2\sin\theta\cos\theta}{r^5}\ep
      \notag\\
    &=\boxed{
       \frac{3\mu_0Ia^2d}{4r^3}
       \sin\theta\cos\theta\,\ep
      },
      \qquad r\gg a,d.
      \label{p5:AB}
\end{align}
Este término no es idénticamente nulo. Su dependencia $1/r^3$
corresponde al término cuadrupolar magnético.

Al calcular $\mathbf B=\nabla\times\mathbf A$, las derivadas
radiales añaden una potencia de $1/r$ y las derivadas angulares
del rotacional también llevan un factor $1/r$. Por tanto,
\[
  \boxed{
    \mathbf A_{\mathrm B}\propto\frac1{r^3},
    \qquad
    \mathbf B_{\mathrm B}\propto\frac1{r^4},
    \qquad r\gg a,d.
  }
\]
Estas potencias describen los términos multipolares dominantes;
sus factores angulares pueden anularse en direcciones particulares.
Las corrientes opuestas cancelan el dipolo y el campo lejano
decrece más rápido que en el caso A.

\end{document}
